%==============================================================================
\section{System Architecture}
\label{sec:architecture}
%==============================================================================

\threatforest{} decomposes threat modeling into six stages, each implemented as an autonomous agent with tool access, orchestrated as a directed graph with verification gates. Figure~\ref{fig:pipeline} shows the full pipeline.

\begin{figure}[t]
\centering
\begin{tikzpicture}[
    node distance=1.2cm and 0.8cm,
    agent/.style={rectangle, draw=blue!60, fill=blue!8, rounded corners, minimum width=2.2cm, minimum height=0.7cm, font=\small},
    verifier/.style={diamond, draw=green!60, fill=green!8, aspect=2, font=\scriptsize, inner sep=1pt},
    arrow/.style={-{Stealth[length=2mm]}, thick},
    retry/.style={-{Stealth[length=2mm]}, thick, dashed, red!60},
]
    \node[agent] (scanner) {Scanner};
    \node[verifier, right=of scanner] (sv) {V};
    \node[agent, right=of sv] (threat) {Threat};
    \node[verifier, right=of threat] (tv) {V};
    \node[agent, below=1.5cm of scanner] (tree) {Tree $\times N$};
    \node[agent, right=of tree] (ttp) {TTP};
    \node[agent, right=of ttp] (mit) {Mitigation};
    \node[verifier, right=of mit] (pv) {V};
    \node[agent, below=1.5cm of tree] (report) {Report};
    \node[verifier, right=of report] (rv) {V};

    \draw[arrow] (scanner) -- (sv);
    \draw[arrow] (sv) -- (threat);
    \draw[arrow] (threat) -- (tv);
    \draw[arrow] (tv) -- node[right, font=\scriptsize] {$\times N$ threats} (tree);
    \draw[arrow] (tree) -- (ttp);
    \draw[arrow] (ttp) -- (mit);
    \draw[arrow] (mit) -- (pv);
    \draw[arrow] (pv) -- (report);
    \draw[arrow] (report) -- (rv);

    \draw[retry] (sv) to[bend right=40] (scanner);
    \draw[retry] (tv) to[bend right=40] (threat);
    \draw[retry] (pv) to[bend left=40] (tree);

    % Parallel bracket
    \draw[decorate, decoration={brace, amplitude=5pt, mirror}, thick] ([yshift=-0.4cm]tree.south west) -- ([yshift=-0.4cm]mit.south east) node[midway, below=0.3cm, font=\scriptsize] {Parallel per threat};
\end{tikzpicture}
\caption{The \threatforest{} pipeline. Boxes are LLM agents; diamonds are deterministic verifiers. Dashed arrows indicate retry edges on verification failure. The tree$\to$TTP$\to$mitigation sub-pipeline runs in parallel for each threat.}
\label{fig:pipeline}
\end{figure}

%----------------------------------------------------------------------
\subsection{Pipeline Overview and Graph Orchestration}
\label{sec:orchestration}
%----------------------------------------------------------------------

The pipeline is implemented as a Strands Graph~\cite{strands2025}---a directed graph where nodes are agent executors and edges carry conditional predicates. Each agent reads its predecessor's output from a shared state directory (\texttt{.threatforest/state/}) and writes structured JSON. This file-based state passing provides three benefits: (1) intermediate outputs are inspectable for debugging, (2) the pipeline is resumable after failures, and (3) parallel agents avoid shared-memory concurrency issues.

The graph defines retry edges: when a verifier detects invalid output, execution flows back to the producing agent for a second attempt, up to a configurable budget of 50 total node executions. This provides self-healing behavior without requiring explicit error handling in each agent.

%----------------------------------------------------------------------
\subsection{Scanner and Threat Generation}
\label{sec:scanner-threat}
%----------------------------------------------------------------------

\textbf{Scanner Agent.} The scanner receives a repository path and uses two tools: a \emph{structural analyzer} that extracts the project's file tree, dependency manifests, and infrastructure-as-code definitions; and a \emph{sandboxed file reader} restricted to the repository directory. The agent produces a \emph{project context} containing: cloud provider, technology stack, services, authentication mechanisms, data flows, and deployment model. A deterministic verifier checks that all required fields are present and non-empty.

\textbf{Threat Agent.} The threat agent receives the project context and generates structured threat statements following the pattern:

\begin{quote}
\emph{A [actor] with [access/knowledge], can [action], which leads to [impact], resulting in [consequence].}
\end{quote}

Each threat includes a title, description, priority, and affected components. The verifier validates JSON structure and ensures each threat references components identified in the scanner context.

%----------------------------------------------------------------------
\subsection{Attack Tree Generation}
\label{sec:tree}
%----------------------------------------------------------------------

Attack trees are generated in parallel: one sub-pipeline per threat, all executing concurrently. Each tree agent receives the scanner context and a single threat, producing a hierarchical attack tree where:

\begin{itemize}
    \item The \textbf{root node} represents the adversary's goal (derived from the threat statement).
    \item \textbf{Intermediate nodes} decompose the goal into sub-goals with AND/OR relationships.
    \item \textbf{Leaf nodes} represent concrete attack techniques.
\end{itemize}

Each step $s_i$ in the tree is a tuple $(id, title, description, parent\_id, is\_leaf, category)$ where $category \in \{\text{fact}, \text{action}, \text{detection}, \text{gate}\}$. The tree structure enables systematic enumeration of attack paths: each root-to-leaf path represents a complete attack scenario.

For a system with $N$ threats, the parallel fan-out produces $N$ independent sub-pipelines (tree $\to$ TTP $\to$ mitigation), providing approximately $N\times$ speedup over sequential execution.

%----------------------------------------------------------------------
\subsection{TTP Mapping: Hybrid Embedding + LLM}
\label{sec:ttp}
%----------------------------------------------------------------------

Mapping attack tree steps to MITRE ATT\&CK techniques is the most technically challenging stage. We employ a two-phase hybrid approach that combines the efficiency of embedding-based retrieval with the contextual reasoning of an LLM.

\textbf{Phase 1: Embedding Retrieval.}
Each attack step description $d_i$ is encoded using ATTACK-BERT~\cite{attackbert}, a sentence-transformer fine-tuned on cybersecurity text. We compute cosine similarity against pre-encoded ATT\&CK technique descriptions $\{t_1, \ldots, t_M\}$ where $M$ is the number of techniques in the ATT\&CK matrix:

\begin{equation}
\text{sim}(d_i, t_j) = \frac{\mathbf{e}(d_i) \cdot \mathbf{e}(t_j)}{\|\mathbf{e}(d_i)\| \, \|\mathbf{e}(t_j)\|}
\label{eq:cosine}
\end{equation}

where $\mathbf{e}(\cdot)$ denotes the ATTACK-BERT embedding function. For each step $d_i$, we retrieve the top-$K$ candidates ranked by similarity:

\begin{equation}
\text{candidates}(d_i) = \text{top-}K\left(\{\text{sim}(d_i, t_j) \mid j = 1, \ldots, M\}\right)
\label{eq:topk}
\end{equation}

We use $K=5$ and a minimum similarity threshold $\tau = 0.3$. The top-1 candidate is selected as the initial mapping.

\textbf{Phase 2: LLM Reviewer.}
The embedding-based top-1 mapping is often correct for well-described attack steps but fails when surface-level textual similarity diverges from semantic intent. For example, a step describing ``credential theft via compromised IAM role'' may match ``T1557.004 Evil Twin'' (a Wi-Fi attack) based on shared vocabulary around ``compromise'' rather than the semantically correct ``T1078.004 Valid Accounts: Cloud Accounts.''

The LLM reviewer agent receives the top-1 summary and has access to a \texttt{get\_ttp\_alternatives} tool that returns the full top-$K$ candidates for any step. The reviewer:
\begin{enumerate}
    \item Examines each top-1 mapping for semantic correctness.
    \item For suspicious mappings, queries alternatives and selects the best match.
    \item Records whether it overrode the embedding's top-1 and provides reasoning.
\end{enumerate}

This hybrid approach leverages the embedding model's efficiency (encoding all steps in a single batch) while using the LLM only for refinement, keeping costs proportional to the number of overrides rather than the total number of steps.

%----------------------------------------------------------------------
\subsection{Mitigation and Verification}
\label{sec:mitigation}
%----------------------------------------------------------------------

\textbf{Mitigation Agent.} The mitigation agent receives the attack trees, TTP mappings, and scanner context, synthesizing per-technique mitigations. Each mitigation includes: the specific AWS service or configuration to implement, the attack steps it addresses, and evidence linking the mitigation to the identified technique. This evidence-based approach ensures mitigations are actionable rather than generic.

\textbf{Verifier Pattern.} Every stage in the pipeline includes a deterministic verifier---a pure function (no LLM) that validates structural properties of the output:

\begin{itemize}
    \item \textbf{Scanner verifier}: All required context fields present and non-empty.
    \item \textbf{Threat verifier}: Valid JSON, each threat has required fields, references valid components.
    \item \textbf{Pipeline verifier}: Attack trees have valid parent-child relationships, TTP mappings reference existing steps, mitigations reference existing techniques.
    \item \textbf{Report verifier}: Output file exists with expected sections.
\end{itemize}

On verification failure, the graph's retry edge routes execution back to the producing agent. This pattern provides quality gates without the cost or non-determinism of LLM-based validation.

\textbf{Report Generator.} The final stage is deterministic (no LLM): it compiles all state files into a structured Markdown report and an interactive HTML dashboard with attack tree visualizations.
